\documentclass[11pt]{article}
\usepackage[a4paper,left=0.85in,right=0.85in,top=0.75in,bottom=0.75in]{geometry}
\usepackage{hyperref}
\hypersetup{hidelinks}
\usepackage{parskip}
\begin{document}
\thispagestyle{empty}
{\Large \textbf{HR Video Script}}\\
Databricks - Systems PhD - Software Engineer\\
Prepared on April 5, 2026

US Style Notes

Start with a smile and steady eye contact.
Use short sentences and pause after each key point.
Keep the tone clear, positive, and direct.
When talking about research, explain the problem, what you built, and why it mattered before naming the venue.

Key Message To Sell Yourself

I can build real systems, explain technical work clearly, and connect research ideas to practical outcomes.

Script

Hi, I am Guanli Liu.

I am based in Melbourne, Australia.

I am applying for the Systems PhD - Software Engineer role at Databricks.

The reason I am applying is that this role is very close to the kind of work I have been doing and the kind of problems I enjoy most. I like database and backend systems work where performance, design, and practical impact all matter.

I bring about six years of research experience in database systems and about three years of engineering experience in backend and data-intensive software.

Recently, I have been working on database systems research at the University of Melbourne. A big part of my work has been physical data design, indexing, benchmarking, and performance evaluation. For example, I built a layout advisory system for data lake style workloads. It supports dataset ingestion, SQL workload parsing, and layout recommendation. I also developed a cost model for physical design trade-offs. This work led to publications at VLDB 2025 and VLDB 2026.

I also built DriftBench to study how system performance changes when workloads or data drift over time. The key point was to measure whether systems stay stable when conditions change, not only whether they look fast in a fixed benchmark. This work was published at VLDB 2026.

Beyond research, I also have practical engineering experience. Earlier in industry, I worked on messaging infrastructure and database performance tuning at Baidu, and I have also built backend data workflows, ingestion pipelines, and practical system tooling in applied engineering settings.

I think I match this role well because the Databricks team cares about database systems, performance, indexes, storage structures, and practical system building. Those are the areas where I have the strongest overlap.

What I can bring is a combination of systems research depth, hands-on engineering, and careful performance thinking. I can build real systems, evaluate them carefully, and explain technical work clearly.

Thank you for your time. I would be excited to speak with you.

\end{document}
